\documentclass[aspectratio=169]{beamer}
\usetheme[showlogo,nosectionpage]{UU}

\usepackage[utf8]{inputenc}
\usepackage[spanish]{babel}
\usepackage{csquotes}
\usepackage{graphicx}
\usepackage{booktabs}
\usepackage{colortbl}
\usepackage{amsmath, amsfonts}
\usepackage{hyperref}
\usepackage{tikz}
\usepackage{fontawesome5}
\usetikzlibrary{positioning, arrows.meta, shapes.geometric, fit, calc}

\definecolor{branchblue}{HTML}{0969DA}
\definecolor{branchgreen}{HTML}{238636}
\definecolor{branchpurple}{HTML}{8250DF}
\definecolor{linkedinblue}{HTML}{0A66C2}

\newcommand{\pitEdition}{9na Edicion}
\newcommand{\pitSession}{Programa completo}

\graphicspath{
  {./}
  {images/}
}

\newcommand{\ProgramCoverSlide}{%
  \begin{frame}[plain]
    \begin{tikzpicture}[remember picture, overlay]
      \fill[OTIred]
        (current page.north west) rectangle
        ([xshift=8mm]current page.south west);
      \shade[left color=OTIred, right color=OTIred!0!white, opacity=0.22]
        ([xshift=8mm]current page.north west) rectangle
        ([xshift=0.38\paperwidth]current page.south west);
    \end{tikzpicture}

    \vspace*{0.45cm}
    \begin{center}
      \includegraphics[height=0.92cm,keepaspectratio]{logos/LOGO-UNI.png}\hspace{0.55cm}
      \includegraphics[height=0.92cm,keepaspectratio]{logos/logo transformacion digital.png}\hspace{0.55cm}
      \includegraphics[height=0.92cm,keepaspectratio]{logos/logo cursos gratuitos.png}
    \end{center}

    \vspace{0.65cm}
    \begin{columns}[c]
      \column{0.64\textwidth}
      \hspace*{0.65cm}
      \begin{minipage}{0.98\linewidth}
        {\usebeamerfont{title}\fontsize{24}{30}\selectfont\bfseries\mbox{Fundamentos de Linux}}\par
        \vspace{0.18cm}
        {\usebeamerfont{subtitle}\color{UUdarkgray}\large \pitSession\ - \pitEdition}\par
        \vspace{0.35cm}
        {\color{OTIred}\rule{3cm}{1.5pt}}\par
        \vspace{0.35cm}
        {\usebeamerfont{author}\textbf{Instructor:} \href{https://www.linkedin.com/in/anthony-gonzalo-quispe-cordova-a39274196}{\color{linkedinblue}\faLinkedin\ Anthony Gonzalo Quispe Cordova}}\par
        \vspace{0.12cm}
        {\usebeamerfont{institute}\color{UUdarkgray}Oficina de Tecnolog\'ias de la Informaci\'on (OTI)}\par
        \vspace{0.12cm}
        {\usebeamerfont{date}\color{UUdarkgray}\today}
      \end{minipage}

      \column{0.28\textwidth}
      \centering
      \begin{tikzpicture}
        \node[rounded corners=8pt, fill=UUlightgray, draw=OTIred!35, minimum width=3.7cm, minimum height=3.0cm, align=center] {
          {\color{green!60!black}\fontsize{42}{46}\selectfont\faLinux}\quad
          {\color{UUdarkblue}\fontsize{42}{46}\selectfont\faTerminal}\\[0.25cm]
          {\large\bfseries Linux + Shell}\\[0.05cm]
          {\scriptsize Base para DevOps y automatizaci\'on}
        };
      \end{tikzpicture}
    \end{columns}
  \end{frame}
}

\newcommand{\MainHeaderSlide}[2]{%
  \begin{frame}[plain]
    \hypertarget{#2}{}
    \vspace*{0.25\textheight}
    {\color{black}\rule{0.28\textwidth}{1.2pt}}\\[0.5em]
    {\usebeamerfont{title}\fontsize{34}{40}\selectfont\bfseries #1}
  \end{frame}
}

\title[Programa Linux]{Fundamentos de Linux}
\subtitle{\pitSession\ - \pitEdition}
\author[Anthony Gonzalo Quispe Cordova]{Instructor: \href{https://www.linkedin.com/in/anthony-gonzalo-quispe-cordova-a39274196}{\color{linkedinblue}\faLinkedin\ Anthony Gonzalo Quispe Cordova}}
\institute[OTI]{Oficina de Tecnolog\'ias de la Informaci\'on (OTI)}
\date{\today}

\begin{document}

\ProgramCoverSlide

% Slides obligatorios segun la plantilla OTI.
\begin{frame}{Flujo de Cursos Orientado a DevOps}
  \begin{center}
    \includegraphics[width=\textwidth,height=0.8\textheight,keepaspectratio]{images/flujo_cursos.png}
  \end{center}
\end{frame}

\begin{frame}{Horarios de los Cursos del Flujo DevOps}
  \begin{center}
  \begin{tikzpicture}[
    card/.style={draw, rounded corners=4pt, minimum width=5.5cm, minimum height=2.6cm, align=center, text width=5.2cm},
    titulo/.style={font=\bfseries\small},
    detalle/.style={font=\footnotesize}
  ]
    \node[card, fill=black!85, draw=yellow, text=white] (c1) at (0,0) {
      {\color{yellow}\faLinux\ \textbf{Fundamentos de Linux}}\\[4pt]
      {\footnotesize \faCalendar\ S\'abados}\\
      {\footnotesize \faClock\ 9:00 a.m. -- 12:00 p.m.}\\[2pt]
      {\scriptsize \faFlag\ 23/05 -- 27/06 \quad | \quad 18 h}
    };
    \node[card, fill=OTIred!15, draw=OTIred!80, right=0.5cm of c1] (c2) {
      {\color{OTIred}\faCodeBranch\ \textbf{Git y GitHub}}\\[4pt]
      {\footnotesize \faCalendar\ Lun, Mar, Jue, Vie}\\
      {\footnotesize \faClock\ 4:00 -- 7:00 p.m.}\\[2pt]
      {\scriptsize \faFlag\ 19/05 -- 28/05 \quad | \quad 18 h}
    };
    \node[card, fill=blue!10, draw=blue!60, below=0.4cm of c1] (c3) {
      {\color{blue!70}\faDocker\ \textbf{Fundamentos de Docker}}\\[4pt]
      {\footnotesize \faCalendar\ Lun, Mar, Vie}\\
      {\footnotesize \faClock\ 4:00 -- 7:00 p.m.}\\[2pt]
      {\scriptsize \faFlag\ 29/05 -- 09/06 \quad | \quad 18 h}
    };
    \node[card, fill=orange!15, draw=orange!70, below=0.4cm of c2] (c4) {
      {\color{orange!80!black}\faServer\ \textbf{Ansible: Automatizaci\'on}}\\[4pt]
      {\footnotesize \faCalendar\ S\'abados}\\
      {\footnotesize \faClock\ 2:00 -- 5:00 p.m.}\\[2pt]
      {\scriptsize \faFlag\ 23/05 -- 27/06 \quad | \quad 18 h}
    };
  \end{tikzpicture}
  \end{center}
\end{frame}

\begin{frame}{\'Indice General}
  \small
  \begin{columns}[t]
    \column{0.48\textwidth}
    \textbf{Linux y Terminal desde Cero}
    \begin{enumerate}
      \item \hyperlink{idx-linux-importancia}{Linux y su importancia actual}
      \item \hyperlink{idx-ubuntu-virtualbox}{Ubuntu Server en VirtualBox}
      \item \hyperlink{idx-terminal-comandos}{Terminal y comandos esenciales}
      \item \hyperlink{idx-laboratorio-practico}{Laboratorio practico}
    \end{enumerate}
    \column{0.48\textwidth}
    \textbf{Usuarios, Permisos y Sistema de Archivos}
    \begin{enumerate}
      \item \hyperlink{idx-identidad-linux}{Identidad: usuarios, grupos y sudo}
      \item \hyperlink{idx-permisos-linux}{Permisos, propietarios y modos}
      \item \hyperlink{idx-editores-config}{Edicion de configuracion}
      \item \hyperlink{idx-lab-multiusuario}{Laboratorio multiusuario}
    \end{enumerate}
  \end{columns}
\end{frame}

\section{Linux y Terminal desde Cero}
\MainHeaderSlide{Linux y Terminal desde Cero}{idx-semana-1}

\subsection{Linux y su importancia actual}
\MainHeaderSlide{Linux y su importancia actual}{idx-linux-importancia}

\subsection{Que es Linux}
\begin{frame}{Que es Linux}
  \begin{columns}[c]
    \column{0.62\textwidth}
    \begin{itemize}
      \item \textbf{Linux} es un sistema operativo libre, de codigo abierto y tipo Unix.
      \item Su componente central es el \textbf{kernel}: administra CPU, memoria, disco, red y dispositivos.
      \item Una \textbf{distribucion} agrega herramientas, instalador, repositorios y configuracion lista para usar.
      \item Ejemplos comunes: Ubuntu, Debian, Fedora, Rocky Linux, Kali Linux y Arch Linux.
    \end{itemize}
    \column{0.33\textwidth}
    \centering
    {\color{OTIred}\fontsize{58}{64}\selectfont\faLinux}\\[0.35cm]
    {\large\textbf{Base de servidores modernos}}
  \end{columns}
\end{frame}

\begin{frame}{Por que Linux domina servidores, nube y ciberseguridad}
  \begin{columns}[t]
    \column{0.33\textwidth}
    \begin{block}{\faServer\ Servidores}
      \small
      Estabilidad, bajo consumo, automatizacion y administracion remota por terminal.
    \end{block}
    \column{0.33\textwidth}
    \begin{exampleblock}{\faCloud\ Nube}
      \small
      La mayoria de servicios cloud, contenedores y pipelines DevOps trabajan sobre Linux.
    \end{exampleblock}
    \column{0.33\textwidth}
    \begin{alertblock}{\faUserShield\ Seguridad}
      \small
      Herramientas de analisis, hardening, monitoreo y respuesta suelen ejecutarse en Linux.
    \end{alertblock}
  \end{columns}
  \vfill
  \begin{beamercolorbox}[rounded=true, shadow=false, sep=4pt, center]{block body}
    \small Aprender Linux es aprender la base operativa de infraestructura, DevOps, redes y seguridad.
  \end{beamercolorbox}
\end{frame}

\begin{frame}{Distribuciones: el mismo nucleo, distintos objetivos}
  \begin{center}
  \begin{tikzpicture}[
    box/.style={draw, rounded corners=6pt, minimum width=4.8cm, minimum height=1.8cm, align=center, text width=4.4cm, font=\small}
  ]
    \node[box, fill=orange!15, draw=orange!70] (b1) at (0,0) {
      {\color{orange!80!black}\faUbuntu\ \textbf{Ubuntu Server}}\\[3pt]
      Servidores y aprendizaje.
    };
    \node[box, fill=green!15, draw=green!60!black] (b2) at (5.5,0) {
      {\color{green!60!black}\faLinux\ \textbf{Debian}}\\[3pt]
      Estabilidad y base comunitaria.
    };
    \node[box, fill=blue!10, draw=blue!60] (b3) at (0,-2.45) {
      {\color{blue!70}\faServer\ \textbf{Rocky / RHEL}}\\[3pt]
      Empresa y produccion.
    };
    \node[box, fill=OTIred!12, draw=OTIred!70] (b4) at (5.5,-2.45) {
      {\color{OTIred}\faUserShield\ \textbf{Kali Linux}}\\[3pt]
      Seguridad ofensiva y laboratorios.
    };
  \end{tikzpicture}
  \end{center}
\end{frame}

\subsection{Distros, escritorios y paquetes}
\begin{frame}{Familias de distribuciones Linux}
  \small
  \begin{columns}[t]
    \column{0.33\textwidth}
    \begin{block}{Debian}
      \begin{itemize}
        \item Debian
        \item Ubuntu
        \item Linux Mint
        \item Kali Linux
      \end{itemize}
      \centering\texttt{apt}, \texttt{dpkg}
    \end{block}
    \column{0.33\textwidth}
    \begin{alertblock}{Red Hat}
      \begin{itemize}
        \item RHEL
        \item Fedora
        \item Rocky Linux
        \item AlmaLinux
      \end{itemize}
      \centering\texttt{dnf}, \texttt{rpm}
    \end{alertblock}
    \column{0.33\textwidth}
    \begin{exampleblock}{Arch / SUSE}
      \begin{itemize}
        \item Arch Linux
        \item Manjaro
        \item openSUSE
        \item SUSE Linux
      \end{itemize}
      \centering\texttt{pacman}, \texttt{zypper}
    \end{exampleblock}
  \end{columns}
  \vfill
  \begin{beamercolorbox}[rounded=true, shadow=false, sep=4pt, center]{block body}
    \small La familia define repositorios, formato de paquetes, comandos de instalacion y ciclo de actualizaciones.
  \end{beamercolorbox}
\end{frame}

\begin{frame}{Entornos de escritorio vs gestores de ventanas}
  \footnotesize
  \begin{columns}[t]
    \column{0.5\textwidth}
    \begin{block}{Entorno de escritorio}
      \begin{itemize}
        \item Experiencia grafica completa.
        \item Paneles, menus, configuracion y apps base.
        \item Ejemplos: GNOME, KDE Plasma, XFCE, Cinnamon.
      \end{itemize}
    \end{block}
    \column{0.5\textwidth}
    \begin{exampleblock}{Gestor de ventanas}
      \begin{itemize}
        \item Controla como se abren y acomodan ventanas.
        \item Puede ser flotante o tipo mosaico.
        \item Ejemplos: i3, Openbox, AwesomeWM, Sway.
      \end{itemize}
    \end{exampleblock}
  \end{columns}
  \vspace{0.12cm}
  \begin{alertblock}{Importante para servidores}
    Ubuntu Server normalmente se administra \textbf{sin entorno grafico}; por eso la terminal es la herramienta principal.
  \end{alertblock}
\end{frame}

\begin{frame}{Paqueterias y gestores de paquetes}
  \small
  \begin{center}
  \begin{tabular}{p{0.23\textwidth} p{0.19\textwidth} p{0.16\textwidth} p{0.31\textwidth}}
    \toprule
    Familia & Formato & Gestor & Ejemplo \\
    \midrule
    Debian / Ubuntu & \texttt{.deb} & \texttt{apt} & \texttt{sudo apt install vim} \\
    Red Hat / Fedora & \texttt{.rpm} & \texttt{dnf} & \texttt{sudo dnf install vim} \\
    Arch / Manjaro & \texttt{pkg.tar.zst} & \texttt{pacman} & \texttt{sudo pacman -S vim} \\
    openSUSE / SUSE & \texttt{.rpm} & \texttt{zypper} & \texttt{sudo zypper install vim} \\
    \bottomrule
  \end{tabular}
  \end{center}
  \vspace{0.25cm}
  \begin{columns}[t]
    \column{0.5\textwidth}
    \begin{block}{Repositorio}
      \small Servidor confiable desde donde la distro descarga software y actualizaciones.
    \end{block}
    \column{0.5\textwidth}
    \begin{exampleblock}{Paquete}
      \small Archivo instalable que incluye programa, metadatos, version y dependencias.
    \end{exampleblock}
  \end{columns}
\end{frame}

\begin{frame}{Formatos universales de instalacion}
  \small
  \begin{columns}[t]
    \column{0.33\textwidth}
    \begin{block}{Snap}
      \begin{itemize}
        \item Usado por Ubuntu.
        \item Apps autocontenidas.
        \item Comando: \texttt{snap}.
      \end{itemize}
    \end{block}
    \column{0.33\textwidth}
    \begin{exampleblock}{Flatpak}
      \begin{itemize}
        \item Popular en escritorio.
        \item Repos como Flathub.
        \item Comando: \texttt{flatpak}.
      \end{itemize}
    \end{exampleblock}
    \column{0.33\textwidth}
    \begin{alertblock}{AppImage}
      \begin{itemize}
        \item Archivo ejecutable.
        \item No requiere instalacion tradicional.
        \item Util para apps portables.
      \end{itemize}
    \end{alertblock}
  \end{columns}
  \vfill
  \centering
  \small En servidores se priorizan repositorios oficiales; en escritorio aparecen mas formatos universales.
\end{frame}

\subsection{Ubuntu Server en VirtualBox}
\MainHeaderSlide{Ubuntu Server en VirtualBox}{idx-ubuntu-virtualbox}

\subsection{Objetivo del entorno}
\begin{frame}{Objetivo del entorno virtual}
  \begin{columns}[c]
    \column{0.55\textwidth}
    \begin{itemize}
      \item Practicar Linux sin modificar el sistema principal.
      \item Crear un servidor real en una maquina virtual.
      \item Aprender a iniciar sesion y trabajar desde la terminal.
      \item Repetir, romper y reconstruir el laboratorio con seguridad.
    \end{itemize}
    \column{0.4\textwidth}
    \centering
    \begin{tikzpicture}[
      box/.style={draw, rounded corners=4pt, minimum width=3.3cm, minimum height=0.9cm, align=center, font=\small},
      arr/.style={-{Stealth[length=2mm]}, thick, OTIred}
    ]
      \node[box, fill=UUlightgray] (host) {\faLaptop\ Equipo fisico};
      \node[box, fill=UUsoftred, below=0.45cm of host] (vbox) {\faCube\ VirtualBox};
      \node[box, fill=UUlightgray, below=0.45cm of vbox] (vm) {\faServer\ Ubuntu Server};
      \node[box, fill=UUsoftred, below=0.45cm of vm] (term) {\faTerminal\ Terminal};
      \draw[arr] (host) -- (vbox);
      \draw[arr] (vbox) -- (vm);
      \draw[arr] (vm) -- (term);
    \end{tikzpicture}
  \end{columns}
\end{frame}

\begin{frame}{Requisitos para la instalacion}
  \small
  \begin{columns}[t]
    \column{0.5\textwidth}
    \begin{block}{Software}
      \begin{itemize}
        \item Oracle VirtualBox instalado.
        \item ISO de Ubuntu Server.
        \item Conexion a internet para descargas.
      \end{itemize}
    \end{block}
    \column{0.5\textwidth}
    \begin{block}{Recursos sugeridos}
      \begin{itemize}
        \item RAM: 2 GB como minimo.
        \item CPU: 2 nucleos si estan disponibles.
        \item Disco: 20 GB dinamico.
        \item Red: NAT para empezar.
      \end{itemize}
    \end{block}
  \end{columns}
  \vfill
  \begin{alertblock}{Antes de iniciar}
    \small Verifica que la virtualizacion este habilitada en BIOS/UEFI si VirtualBox no permite crear maquinas de 64 bits.
  \end{alertblock}
\end{frame}

\subsection{Pasos de instalacion}
\begin{frame}{Crear la maquina virtual}
  \begin{enumerate}
    \item Crear una nueva VM: \textbf{Linux} / \textbf{Ubuntu (64-bit)}.
    \item Asignar memoria y procesadores.
    \item Crear disco duro virtual VDI, reservado dinamicamente.
    \item Montar la ISO de Ubuntu Server en la unidad optica.
    \item Iniciar la VM y seguir el instalador.
  \end{enumerate}
  \vfill
  \begin{beamercolorbox}[rounded=true, shadow=false, sep=4pt, center]{block body}
    \small Meta: terminar con un servidor Ubuntu arrancando y pidiendo usuario/contrasena en consola.
  \end{beamercolorbox}
\end{frame}

\begin{frame}{Decisiones dentro del instalador}
  \begin{columns}[t]
    \column{0.5\textwidth}
    \begin{block}{Configuracion base}
      \begin{itemize}
        \item Idioma y teclado.
        \item Red por DHCP.
        \item Disco completo para la VM.
        \item Usuario administrador del laboratorio.
      \end{itemize}
    \end{block}
    \column{0.5\textwidth}
    \begin{exampleblock}{Perfil recomendado}
      \begin{itemize}
        \item Nombre servidor: \texttt{linux-lab}.
        \item Usuario: \texttt{pit}.
        \item Instalar OpenSSH si se usara acceso remoto.
        \item No instalar extras innecesarios.
      \end{itemize}
    \end{exampleblock}
  \end{columns}
\end{frame}

\begin{frame}[fragile]{Primer ingreso al servidor}
  \begin{itemize}
    \item Inicia sesion con el usuario creado durante la instalacion.
    \item Confirma que estas en Linux y que tienes conexion.
    \item Ubica tu carpeta personal y reconoce el prompt.
  \end{itemize}
  \vspace{0.25cm}
  \begin{exampleblock}{Comandos iniciales}
\begin{verbatim}
whoami
hostname
pwd
ip addr
\end{verbatim}
  \end{exampleblock}
\end{frame}

\subsection{Terminal y comandos esenciales}
\MainHeaderSlide{Terminal y comandos esenciales}{idx-terminal-comandos}

\subsection{Terminal, shell y rutas}
\begin{frame}{Terminal, shell y sistema de archivos}
  \begin{itemize}
    \item \textbf{Terminal:} interfaz donde escribimos comandos.
    \item \textbf{Shell:} programa que interpreta esos comandos; en Ubuntu suele ser Bash.
    \item \textbf{Ruta absoluta:} empieza desde la raiz, por ejemplo \texttt{/home/pit}.
    \item \textbf{Ruta relativa:} parte desde la ubicacion actual, por ejemplo \texttt{Documentos/notas.txt}.
  \end{itemize}
  \vfill
  \begin{beamercolorbox}[rounded=true, shadow=false, sep=4pt, center]{block body}
    \small En Linux todo vive dentro de un arbol que empieza en \texttt{/}; no existen unidades \texttt{C:} o \texttt{D:}.
  \end{beamercolorbox}
\end{frame}

\begin{frame}{Estructura basica del sistema}
  \begin{center}
  \footnotesize
  \begin{tabular}{>{\ttfamily}l p{0.68\textwidth}}
    \toprule
    \normalfont Directorio & Uso principal \\
    \midrule
    /        & Raiz del sistema \\
    /home    & Carpetas de usuarios \\
    /etc     & Configuracion del sistema \\
    /var     & Datos variables como logs \\
    /tmp     & Archivos temporales \\
    /bin     & Comandos esenciales \\
    \bottomrule
  \end{tabular}
  \end{center}
  \vspace{0.35cm}
  \centering
  \small La primera practica se concentrara en moverse por este arbol y crear estructura propia.
\end{frame}

\subsection{Navegar, crear y gestionar archivos}
\begin{frame}[fragile]{Comandos para navegar}
  \begin{columns}[t]
    \column{0.5\textwidth}
    \begin{block}{Ubicacion}
\begin{verbatim}
pwd
ls
ls -la
\end{verbatim}
    \end{block}
    \column{0.5\textwidth}
    \begin{block}{Movimiento}
\begin{verbatim}
cd carpeta
cd ..
cd ~
clear
\end{verbatim}
    \end{block}
  \end{columns}
\end{frame}

\begin{frame}[fragile]{Comandos para crear y gestionar}
  \small
  \begin{columns}[t]
    \column{0.5\textwidth}
    \begin{block}{Crear}
\begin{verbatim}
mkdir proyecto
mkdir -p app/src
touch notas.txt
nano notas.txt
\end{verbatim}
    \end{block}
    \column{0.5\textwidth}
    \begin{block}{Gestionar}
\begin{verbatim}
cp origen destino
mv origen destino
rm archivo
rm -r carpeta
\end{verbatim}
    \end{block}
  \end{columns}
  \vfill
  \begin{alertblock}{Precaucion}
    \small Antes de eliminar con \texttt{rm}, confirma tu ruta con \texttt{pwd} y revisa con \texttt{ls}.
  \end{alertblock}
\end{frame}

\begin{frame}[fragile]{Ver contenido de archivos}
  \begin{columns}[t]
    \column{0.5\textwidth}
    \begin{block}{Lectura rapida}
\begin{verbatim}
cat archivo
head archivo
tail archivo
less archivo
\end{verbatim}
    \end{block}
    \column{0.5\textwidth}
    \begin{exampleblock}{Buenas practicas}
      \begin{itemize}
        \item Usar nombres claros.
        \item Evitar espacios en rutas al empezar.
        \item Organizar por carpetas.
        \item Leer mensajes de error con calma.
      \end{itemize}
    \end{exampleblock}
  \end{columns}
\end{frame}

\subsection{Laboratorio practico}
\MainHeaderSlide{Laboratorio practico}{idx-laboratorio-practico}

\subsection{Construir una estructura real}
\begin{frame}[fragile]{Laboratorio: estructura de proyecto en Linux}
  \begin{columns}[c]
    \column{0.55\textwidth}
    \begin{itemize}
      \item Crear una carpeta principal del curso.
      \item Separar documentacion, scripts, datos y respaldos.
      \item Crear archivos iniciales.
      \item Copiar, mover y renombrar elementos.
      \item Validar la estructura final desde terminal.
    \end{itemize}
    \column{0.4\textwidth}
    \footnotesize
    \begin{block}{Estructura esperada}
\begin{verbatim}
proyecto-linux/
  docs/
  scripts/
  data/
  backups/
  README.txt
\end{verbatim}
    \end{block}
  \end{columns}
\end{frame}

\begin{frame}[fragile]{Paso 1: crear carpetas y archivos}
  \begin{exampleblock}{Comandos}
\begin{verbatim}
cd ~
mkdir -p proyecto-linux/{docs,scripts,data,backups}
cd proyecto-linux
touch README.txt docs/notas.txt scripts/inicio.sh
pwd
ls -la
\end{verbatim}
  \end{exampleblock}
  \vfill
  \small La opcion \texttt{-p} permite crear carpetas intermedias sin error si ya existen.
\end{frame}

\begin{frame}[fragile]{Paso 2: escribir y revisar contenido}
  \begin{exampleblock}{Comandos}
\begin{verbatim}
echo "Proyecto Linux - Bloque inicial" > README.txt
echo "Comandos practicados:" > docs/notas.txt
echo "pwd, ls, cd, mkdir, touch" >> docs/notas.txt
cat README.txt
cat docs/notas.txt
\end{verbatim}
  \end{exampleblock}
  \vfill
  \small En esta semana solo usaremos escritura basica para documentar el laboratorio.
\end{frame}

\begin{frame}[fragile]{Paso 3: copiar, mover y validar}
  \begin{exampleblock}{Comandos}
\begin{verbatim}
cp README.txt backups/README.bak
mv docs/notas.txt docs/comandos-basicos.txt
ls -R
\end{verbatim}
  \end{exampleblock}
  \vfill
  \begin{beamercolorbox}[rounded=true, shadow=false, sep=4pt, center]{block body}
    \small Resultado: un proyecto ordenado, navegable y documentado desde la terminal.
  \end{beamercolorbox}
\end{frame}

\subsection{Cierre}
\begin{frame}{Checklist de aprendizaje}
  \begin{itemize}
    \item Puedo explicar que es Linux y por que se usa en servidores.
    \item Puedo diferenciar distribuciones, escritorios y gestores de paquetes.
    \item Puedo describir como se instala Ubuntu Server en VirtualBox.
    \item Puedo moverme por el sistema de archivos desde la terminal.
    \item Puedo crear, copiar, mover y revisar archivos con comandos basicos.
  \end{itemize}
\end{frame}

\begin{frame}{Cierre: Linux y Terminal}
  \begin{enumerate}
    \item Linux es la base de servidores, nube, DevOps y ciberseguridad.
    \item Ubuntu Server en VirtualBox permite practicar en un entorno real y controlado.
    \item La terminal y Bash permiten controlar el sistema mediante comandos.
    \item El sistema de archivos parte desde \texttt{/} y organiza todo en carpetas.
    \item El laboratorio construyo una estructura de proyecto real en Linux.
  \end{enumerate}
  \vspace{0.3cm}
  \begin{alertblock}{Siguiente bloque}
    \textbf{Tema siguiente:} usuarios, permisos y sistema de archivos. Veremos como Linux organiza el acceso y protege recursos del sistema.
  \end{alertblock}
\end{frame}

\section{Usuarios, Permisos y Sistema de Archivos}

\subsection{Identidad en Linux}
\MainHeaderSlide{Identidad: usuarios, grupos y sudo}{idx-identidad-linux}

\begin{frame}{Objetivos: usuarios, permisos y archivos}
  \begin{columns}[c]
    \column{0.58\textwidth}
    \begin{itemize}
      \item Entender como Linux controla el acceso mediante usuarios, grupos y permisos.
      \item Leer permisos y propietarios sin depender de memoria mecanica.
      \item Usar \texttt{sudo}, \texttt{chmod}, \texttt{chown} y \texttt{chgrp} con criterio.
      \item Editar archivos de configuracion con Nano y una base practica de Vim.
      \item Construir un entorno multiusuario con accesos protegidos por rol.
    \end{itemize}
    \column{0.36\textwidth}
    \centering
    {\color{OTIred}\fontsize{52}{58}\selectfont\faUserShield}\\[0.25cm]
    {\large\textbf{Seguridad por identidad, propiedad y permisos}}
  \end{columns}
\end{frame}

\begin{frame}{Como Linux decide quien puede hacer que}
  \small
  \begin{columns}[t]
    \column{0.33\textwidth}
    \begin{block}{Usuario}
      Identidad que ejecuta procesos, crea archivos y tiene una carpeta de trabajo.
    \end{block}
    \column{0.33\textwidth}
    \begin{exampleblock}{Grupo}
      Conjunto de usuarios usado para compartir permisos sobre recursos.
    \end{exampleblock}
    \column{0.33\textwidth}
    \begin{alertblock}{Otros}
      Cualquier usuario que no sea propietario ni pertenezca al grupo del recurso.
    \end{alertblock}
  \end{columns}
  \vfill
  \begin{beamercolorbox}[rounded=true, shadow=false, sep=4pt, center]{block body}
    \small Cada archivo tiene propietario, grupo propietario y permisos separados para usuario, grupo y otros.
  \end{beamercolorbox}
\end{frame}

\begin{frame}[fragile]{Usuarios, root y sudo}
  \small
  \begin{columns}[t]
    \column{0.5\textwidth}
    \begin{block}{Identidades clave}
      \begin{itemize}
        \item \textbf{root:} administrador total del sistema.
        \item \textbf{Usuario regular:} trabaja con permisos limitados.
        \item \textbf{sudo:} eleva privilegios solo para una accion concreta.
      \end{itemize}
    \end{block}
    \column{0.5\textwidth}
    \begin{exampleblock}{Comandos de reconocimiento}
\begin{verbatim}
whoami
id
groups
sudo -l
\end{verbatim}
    \end{exampleblock}
  \end{columns}
  \vfill
  \begin{alertblock}{Regla operativa}
    Operar como usuario normal y elevar privilegios solo cuando sea necesario.
  \end{alertblock}
\end{frame}

\begin{frame}[fragile]{Gestion basica de cuentas y grupos}
  \small
  \begin{columns}[t]
    \column{0.5\textwidth}
    \begin{block}{Usuarios}
\begin{verbatim}
sudo adduser ana
sudo passwd ana
id ana
sudo deluser ana
\end{verbatim}
    \end{block}
    \column{0.5\textwidth}
    \begin{block}{Grupos}
\begin{verbatim}
sudo groupadd devops
sudo usermod -aG devops ana
groups ana
sudo groupdel devops
\end{verbatim}
    \end{block}
  \end{columns}
  \vfill
  \small Usar \texttt{usermod -aG}: \texttt{-a} agrega sin borrar otros grupos; \texttt{-G} define grupos suplementarios.
\end{frame}

\begin{frame}{Archivos del sistema relacionados con identidad}
  \footnotesize
  \begin{center}
  \begin{tabular}{>{\ttfamily}p{0.22\textwidth} p{0.62\textwidth}}
    \toprule
    \normalfont Archivo & Funcion principal \\
    \midrule
    /etc/passwd & Lista usuarios, UID, GID principal, carpeta personal y shell. \\
    /etc/shadow & Guarda hashes de contrasenas y politicas de expiracion. \\
    /etc/group & Define grupos y miembros suplementarios. \\
    /etc/sudoers & Controla quien puede ejecutar comandos con privilegios. \\
    \bottomrule
  \end{tabular}
  \end{center}
  \vspace{0.25cm}
  \begin{alertblock}{Cuidado}
    \small No editar \texttt{/etc/sudoers} con un editor comun. Usar \texttt{visudo} para validar sintaxis antes de guardar.
  \end{alertblock}
\end{frame}

\subsection{Permisos y propietarios}
\MainHeaderSlide{Permisos, propietarios y modos}{idx-permisos-linux}

\begin{frame}[fragile]{Leer permisos con ls -l}
  \begin{exampleblock}{Salida tipica}
\begin{verbatim}
-rw-r----- 1 ana devops 1240 may 30 14:10 reporte.txt
drwxr-x--- 2 root devops 4096 may 30 14:15 proyectos
\end{verbatim}
  \end{exampleblock}
  \small
  \begin{columns}[t]
    \column{0.5\textwidth}
    \begin{itemize}
      \item Primer caracter: tipo de archivo.
      \item Luego vienen permisos de usuario, grupo y otros.
    \end{itemize}
    \column{0.5\textwidth}
    \begin{itemize}
      \item Despues aparecen propietario y grupo.
      \item Fechas, tamanos y nombres ayudan a validar cambios.
    \end{itemize}
  \end{columns}
\end{frame}

\begin{frame}{Permisos rwx}
  \small
  \begin{columns}[t]
    \column{0.33\textwidth}
    \begin{block}{\texttt{r} lectura}
      Permite leer un archivo o listar nombres dentro de un directorio.
    \end{block}
    \column{0.33\textwidth}
    \begin{exampleblock}{\texttt{w} escritura}
      Permite modificar un archivo o crear, borrar y renombrar dentro de un directorio.
    \end{exampleblock}
    \column{0.33\textwidth}
    \begin{alertblock}{\texttt{x} ejecucion}
      Permite ejecutar un archivo o atravesar un directorio para acceder a su contenido.
    \end{alertblock}
  \end{columns}
  \vfill
  \centering
  \small En directorios, \texttt{x} significa acceso/travesia; sin \texttt{x}, listar no basta para entrar.
\end{frame}

\begin{frame}{Modo simbolico y modo octal}
  \small
  \begin{columns}[t]
    \column{0.5\textwidth}
    \begin{block}{Simbolico}
      \begin{itemize}
        \item \texttt{u}: usuario propietario.
        \item \texttt{g}: grupo propietario.
        \item \texttt{o}: otros.
        \item \texttt{a}: todos.
        \item Operadores: \texttt{+}, \texttt{-}, \texttt{=}.
      \end{itemize}
    \end{block}
    \column{0.5\textwidth}
    \begin{exampleblock}{Octal}
      \begin{itemize}
        \item \texttt{r = 4}
        \item \texttt{w = 2}
        \item \texttt{x = 1}
        \item \texttt{755 = rwx r-x r-x}
        \item \texttt{640 = rw- r-- ---}
      \end{itemize}
    \end{exampleblock}
  \end{columns}
\end{frame}

\begin{frame}[fragile]{chmod: cambiar permisos}
  \begin{columns}[t]
    \column{0.5\textwidth}
    \begin{block}{Modo simbolico}
\begin{verbatim}
chmod u+x script.sh
chmod g+w reporte.txt
chmod o-r secreto.txt
chmod a=r archivo.txt
\end{verbatim}
    \end{block}
    \column{0.5\textwidth}
    \begin{block}{Modo octal}
\begin{verbatim}
chmod 755 script.sh
chmod 644 index.html
chmod 640 reporte.txt
chmod 700 carpeta_privada
\end{verbatim}
    \end{block}
  \end{columns}
  \vfill
  \begin{alertblock}{Evitar}
    \small \texttt{chmod 777} no debe ser la solucion por defecto: abre lectura, escritura y ejecucion a todos.
  \end{alertblock}
\end{frame}

\begin{frame}[fragile]{chown y chgrp: cambiar propiedad}
  \small
  \begin{columns}[t]
    \column{0.5\textwidth}
    \begin{block}{Propietario y grupo}
\begin{verbatim}
sudo chown ana reporte.txt
sudo chown ana:devops reporte.txt
sudo chown -R ana:devops proyecto/
\end{verbatim}
    \end{block}
    \column{0.5\textwidth}
    \begin{block}{Solo grupo}
\begin{verbatim}
sudo chgrp devops reporte.txt
sudo chgrp -R devops proyecto/
ls -l
\end{verbatim}
    \end{block}
  \end{columns}
  \vfill
  \small Cambiar propiedad requiere privilegios porque modifica quien controla el recurso.
\end{frame}

\begin{frame}{Permisos recomendados en situaciones comunes}
  \footnotesize
  \begin{center}
  \begin{tabular}{p{0.28\textwidth} p{0.18\textwidth} p{0.38\textwidth}}
    \toprule
    Caso & Permiso & Criterio \\
    \midrule
    Archivo publico de lectura & \texttt{644} & Propietario escribe; todos leen. \\
    Script ejecutable & \texttt{755} & Propietario modifica; otros ejecutan. \\
    Archivo sensible & \texttt{600} & Solo propietario lee y escribe. \\
    Carpeta privada & \texttt{700} & Solo propietario entra y modifica. \\
    Carpeta compartida por grupo & \texttt{770} & Usuario y grupo trabajan; otros no entran. \\
    \bottomrule
  \end{tabular}
  \end{center}
  \vspace{0.25cm}
  \centering
  \small El permiso correcto depende del rol del archivo, no de una receta unica.
\end{frame}

\begin{frame}{Enlaces simbolicos y enlaces duros}
  \small
  \begin{columns}[t]
    \column{0.5\textwidth}
    \begin{block}{Enlace simbolico}
      Apunta a una ruta. Si el destino cambia o se borra, el enlace puede quedar roto.
      \vspace{0.15cm}

      \texttt{ln -s destino enlace}
    \end{block}
    \column{0.5\textwidth}
    \begin{exampleblock}{Enlace duro}
      Comparte el mismo inodo del archivo. No es una copia separada.
      \vspace{0.15cm}

      \texttt{ln destino enlace}
    \end{exampleblock}
  \end{columns}
  \vfill
  \small Para administracion diaria y configuraciones, el enlace simbolico suele ser el mas visible y facil de interpretar.
\end{frame}

\subsection{Editar configuracion}
\MainHeaderSlide{Edicion de configuracion}{idx-editores-config}

\begin{frame}[fragile]{Nano: edicion directa para empezar}
  \begin{columns}[t]
    \column{0.48\textwidth}
    \begin{itemize}
      \item Simple para editar archivos pequenos.
      \item Muestra atajos en la parte inferior.
      \item Es util cuando el objetivo es corregir rapido y guardar.
    \end{itemize}
    \column{0.52\textwidth}
    \begin{exampleblock}{Flujo basico}
\begin{verbatim}
nano archivo.conf
Ctrl + O   guardar
Enter      confirmar nombre
Ctrl + X   salir
\end{verbatim}
    \end{exampleblock}
  \end{columns}
\end{frame}

\begin{frame}[fragile]{Vim basico: sobrevivir en servidores}
  \small
  \begin{columns}[t]
    \column{0.5\textwidth}
    \begin{block}{Modos esenciales}
\begin{verbatim}
vim archivo.conf
i       insertar
Esc     modo normal
:w      guardar
:q      salir
:wq     guardar y salir
:q!     salir sin guardar
\end{verbatim}
    \end{block}
    \column{0.5\textwidth}
    \begin{exampleblock}{Acciones utiles}
\begin{verbatim}
/texto  buscar
dd      borrar linea
u       deshacer
yy      copiar linea
p       pegar
\end{verbatim}
    \end{exampleblock}
  \end{columns}
\end{frame}

\begin{frame}{Buenas practicas al editar configuracion}
  \begin{enumerate}
    \item Identificar el archivo exacto y confirmar ruta con \texttt{pwd} o ruta absoluta.
    \item Hacer copia antes de cambiar: \texttt{sudo cp archivo archivo.bak}.
    \item Editar el minimo necesario y guardar con comentarios claros si aplica.
    \item Validar sintaxis cuando exista comando de prueba.
    \item Reiniciar o recargar el servicio solo despues de validar.
  \end{enumerate}
  \vfill
  \begin{alertblock}{Criterio}
    \small En servidores, editar sin copia y sin validacion convierte un cambio pequeno en una caida innecesaria.
  \end{alertblock}
\end{frame}

\subsection{Laboratorio multiusuario}
\MainHeaderSlide{Laboratorio multiusuario}{idx-lab-multiusuario}

\begin{frame}{Laboratorio: entorno multiusuario por rol}
  \begin{columns}[c]
    \column{0.58\textwidth}
    \begin{itemize}
      \item Crear usuarios para simular un equipo de trabajo.
      \item Crear grupos por rol.
      \item Preparar una carpeta compartida.
      \item Asignar propietario, grupo y permisos.
      \item Validar que cada usuario solo acceda a lo que corresponde.
    \end{itemize}
    \column{0.36\textwidth}
    \begin{block}{Escenario}
      \small Equipo: \texttt{devops}\\
      Usuarios: \texttt{ana}, \texttt{luis}\\
      Carpeta: \texttt{/srv/proyecto}\\
      Acceso: grupo puede trabajar; otros no.
    \end{block}
  \end{columns}
\end{frame}

\begin{frame}[fragile]{Paso 1: crear usuarios, grupo y carpeta}
  \begin{exampleblock}{Preparar}
\begin{verbatim}
sudo adduser ana
sudo adduser luis
sudo groupadd devops
sudo usermod -aG devops ana
sudo usermod -aG devops luis
sudo mkdir -p /srv/proyecto
\end{verbatim}
  \end{exampleblock}
  \vfill
  \small En un laboratorio real se pueden usar contrasenas temporales y luego eliminarlas al finalizar.
\end{frame}

\begin{frame}[fragile]{Paso 2: asignar propiedad y permisos}
  \begin{exampleblock}{Configurar}
\begin{verbatim}
sudo chown root:devops /srv/proyecto
sudo chmod 770 /srv/proyecto
ls -ld /srv/proyecto
\end{verbatim}
  \end{exampleblock}
  \vfill
  \begin{block}{Interpretacion}
    \small \texttt{770} permite acceso total al propietario y al grupo; bloquea completamente a otros usuarios.
  \end{block}
\end{frame}

\begin{frame}[fragile]{Paso 3: validar accesos}
  \begin{columns}[t]
    \column{0.5\textwidth}
    \begin{block}{Usuario autorizado}
\begin{verbatim}
su - ana
cd /srv/proyecto
touch ana.txt
ls -l
exit
\end{verbatim}
    \end{block}
    \column{0.5\textwidth}
    \begin{block}{Revision como administrador}
\begin{verbatim}
sudo ls -la /srv/proyecto
sudo -u luis bash
touch /srv/proyecto/luis.txt
exit
sudo ls -la /srv/proyecto
\end{verbatim}
    \end{block}
  \end{columns}
\end{frame}

\begin{frame}[fragile]{Paso 4: limpiar el laboratorio}
  \begin{exampleblock}{Limpieza}
\begin{verbatim}
sudo rm -rf /srv/proyecto
sudo deluser ana
sudo deluser luis
sudo groupdel devops
\end{verbatim}
  \end{exampleblock}
  \vfill
  \begin{alertblock}{Antes de borrar}
    \small Confirmar la ruta exacta con \texttt{ls -ld /srv/proyecto}. Nunca ejecutar \texttt{rm -rf} sobre una ruta dudosa.
  \end{alertblock}
\end{frame}

\begin{frame}{Checklist de aprendizaje}
  \begin{itemize}
    \item Puedo explicar la diferencia entre usuario, grupo y otros.
    \item Puedo leer permisos desde una salida \texttt{ls -l}.
    \item Puedo aplicar \texttt{chmod} en modo simbolico y octal.
    \item Puedo cambiar propietario y grupo con \texttt{chown} y \texttt{chgrp}.
    \item Puedo editar configuraciones basicas con Nano y salir correctamente de Vim.
    \item Puedo construir una carpeta compartida protegida por rol.
  \end{itemize}
\end{frame}

\begin{frame}{Resumen: usuarios, permisos y archivos}
  \begin{enumerate}
    \item Linux separa identidad, propiedad y permisos para proteger recursos.
    \item \texttt{sudo} permite privilegios temporales sin trabajar siempre como \texttt{root}.
    \item \texttt{chmod} define accesos; \texttt{chown} y \texttt{chgrp} definen control.
    \item Nano y Vim permiten modificar configuraciones desde servidores sin entorno grafico.
    \item El laboratorio consolida un entorno multiusuario con acceso por rol.
  \end{enumerate}
  \vspace{0.25cm}
  \begin{alertblock}{Siguiente bloque}
    \textbf{Tema siguiente:} procesos, servicios y paquetes. Veremos como monitorear el sistema, controlar procesos y administrar software.
  \end{alertblock}
\end{frame}

\end{document}
